\documentclass[12pt,a4paper]{article}
\usepackage[margin=1in]{geometry}
\usepackage{amsmath,amssymb,amsthm}
\usepackage{hyperref}
\usepackage{enumitem}
\usepackage{booktabs}
\usepackage{array}

% xr-hyper cross-paper references (must follow hyperref).
\usepackage{xr-hyper}
\externaldocument[paper0:]{cascade-series-part0}
\externaldocument[paper1:]{cascade-series-part1}
\externaldocument[paper2:]{cascade-series-part2}
\externaldocument[paper3:]{cascade-series-part3}
\externaldocument[paper4a:]{cascade-series-part4a}
\externaldocument[paper4b:]{cascade-series-part4b}
\externaldocument[paper6:]{cascade-series-part6}
\newcommand{\cascadebase}{https://dickie81.github.io/r-infinite}

\newtheorem{theorem}{Theorem}[section]
\newtheorem{corollary}[theorem]{Corollary}
\newtheorem{lemma}[theorem]{Lemma}
\newtheorem{definition}[theorem]{Definition}
\newtheorem*{remark}{Remark}

\title{\textbf{The Cascade Series --- Part II\,=\,III}\\[6pt]
Quantum Gravity without Quantising Gravity:\\
Why the Quantum, Gravitational, and Matter-Dynamics\\
Projections of the Cascade Are the Same Theorem}
\author{RTAC}
\date{March 2026}

\begin{document}
\setlength{\emergencystretch}{3em}
\maketitle
\phantomsection
\label{paperdoc}

\begin{abstract}
The cascade series tests one hypothesis: the infinite-dimensional unit ball,
descended to four dimensions, is indistinguishable from our universe.
Paper~II~\cite{paper2} derives quantum mechanics as the measurement theory of
a 4D observer in the cascade geometry. Paper~III~\cite{paper3} derives general
relativity as the unique gravitational equation available to that observer.
This paper proves they are the same theorem.

Both projections map the same cascade state on the same unit sphere using the
same propagator under the same identification hypothesis. The propagator
$K(t) = |K|\,e^{-iHt}$ is simultaneously the quantum time-evolution operator
(unitary, probability-preserving) and the source of Lorentzian signature
(oscillatory, $g_{tt} < 0$). These are not two identifications of one
expression---they are one physical fact in two vocabularies. At $d = 4$,
Gleason's theorem forces the Born rule as the unique probability measure and
Lovelock's theorem forces the Einstein equation as the unique gravitational
equation. Two uniquely forced theories from a single source admit no mutual
contradiction: the space of possible conflicts is empty.

The same architecture extends to matter dynamics. Wigner's classification of
relativistic wave equations applied to the cascade Hilbert space (the
irreducibly-complex spinor representation of $\mathrm{Spin}(1,3)$) plus the
cascade-derived mass and gauge structure forces the Dirac equation
$(i\gamma^\mu D_\mu - m)\psi = 0$ at $d=4$. Triple uniqueness from a single
source: Gleason for QM, Lovelock for GR, Wigner for matter dynamics. The
cascade-native gyromagnetic ratio $g = \chi(S^{d-1}) = 2$ follows as a
corollary, with the Schwinger 1-loop coefficient $1/(2\pi) = 1/(N(0)\cdot
\Gamma(\tfrac{1}{2})^2)$ identified as a cascade primitive ratio.

The standard conflicts between quantum mechanics and general relativity---the
UV catastrophe, the problem of time, background independence, the information
paradox, and non-renormalisability---are artefacts of treating the 4D
projections as independent theories. In the cascade, they share a common
origin. There is nothing to reconcile because there was never a separation.

As a direct application, we derive Bekenstein--Hawking entropy from boundary
dominance alone: a horizon of area $A$ at dimension $d$ hides geometric content
$A/d$. At $d = 4$: $S = A/4$. The factor $1/4$ is the boundary-to-volume
ratio at the observer's dimension. Combining $S = A/4$ with the
Schwarzschild area--mass relation $A = 16\pi M^2$ (a mathematical consequence
of the cascade's own Einstein equation via Birkhoff's theorem) gives the
Hawking temperature $T = 1/(8\pi M)$ by calculus---no semiclassical gravity,
no QFT on curved spacetime, no Bogoliubov transformations.
\end{abstract}

\tableofcontents
\newpage

\section{The Schism and Its Dissolution}

\subsection{The standard conflict}

Quantum mechanics requires a fixed background. General relativity makes the
background dynamical. Combining them produces perturbative
non-renormalisability~\cite{goroff}. Eighty years of effort to resolve this
conflict have produced candidate frameworks but no unique, predictive,
empirically confirmed quantum theory of gravity.

\subsection{Why the cascade dissolves the problem}

The cascade does not quantise gravity. It does not need to. The conflict
arises because quantum mechanics and general relativity are formulated as
independent theories, each claiming to be fundamental. In the cascade,
neither is fundamental. Both are 4D projections of a single 217-dimensional
geometry. They cannot conflict for the same reason that the shadow of a cube
on two different walls cannot be inconsistent: both shadows are determined by
the cube.

\section{Single-Source Structure}

\subsection{Shared ingredients}

Papers~II and~III derive their results from the same inputs.

\begin{center}
\begin{tabular}{lll}
\hline
Ingredient & QM projection (II) & GR projection (III) \\
\hline
State space & Unit sphere $S^{d-1}$ & Unit sphere $S^{d-1}$ \\
Fundamental constant & $\sqrt{\pi} = \Gamma(1/2)$ & $\sqrt{\pi} = \Gamma(1/2)$ \\
Propagator & $K(t) = |K|\,e^{-iHt}$ & $K(t) = |K|\,e^{-iHt}$ \\
Complex structure & $J^2 = -\mathrm{Id}$; $U(1) = \langle J\rangle$
                    forced gauge & $J^2 = -\mathrm{Id}$ \\
Time & Slicing $=$ dimensional collapse & Slicing $=$ dimensional collapse \\
Identification & Def.\ 2.1 of III & Def.\ 2.1 of III \\
\hline
\end{tabular}
\end{center}

No entry differs between columns. The two projections draw on the same state
space, the same constant, the same propagator, the same complex structure, the
same time, and the same hypothesis.

\begin{theorem}[Single-source structure]\label{thm:single}
The quantum and gravitational projections are functions of the same cascade
state $|\Psi\rangle \in S^{d-1}$. Neither projection introduces any structure
not present in the cascade geometry.
\end{theorem}

\begin{proof}
Paper~II derives the 4D Hilbert space, the complex Born rule
$|\langle u, v\rangle_\mathbb{C}|^2$ on the $U(1)$-quotient projective
state space $\mathbb{C}P^{n-1}$ (with $U(1)$ forced by Theorem~7.10 of
Paper~II, $\mathit{thm{:}U1{\text{-}}gauge{\text{-}}forced}$), and unitary
evolution from orthogonality, concentration of measure, and the forced
precession---all properties of $S^{d-1}$. Paper~III derives the 4D
metric, Einstein equation, and Lorentzian signature from the foliation
metric, Lovelock's theorem, and the same forced precession---all
properties of $S^{d-1}$. No external structure is invoked in either
derivation. Both are read from the same sphere.
\end{proof}

\begin{remark}[The identification hypothesis is shared]
Definition~2.1 of Paper~III does not say ``the cascade is quantum
mechanical'' or ``the cascade is gravitational.'' It says ``the cascade's
geometry is our physics.'' The quantum and gravitational content are
consequences, not inputs.
\end{remark}

\subsection{The state space is the Hilbert space is the configuration space}

In standard physics, the Hilbert space of QM and the configuration space of GR
are separate mathematical objects. In the cascade, they are the same object:
the unit sphere $S^{d-1}$.

\begin{theorem}[State space identity]\label{thm:statespace}
The boundary $S^{d-1}$ that carries the quantum state (Paper~II) is the
same $S^{d-1}$ whose geometry determines the metric (Paper~III). A cascade
state $|\Psi\rangle \in S^{d-1}$ simultaneously specifies the quantum
amplitudes and the geometric configuration. There is one sphere, not two
sets of degrees of freedom.
\end{theorem}

\begin{proof}
Paper~II: the state space is $S^{d-1}$ via boundary dominance
($\Omega_{d-1}/V_d = d$; the cascade's content is its boundary).
Paper~III: the metric is the foliation of $B^{d+1}$ into cross-sections,
with the induced geometry on $S^{d-1}$ determining the 4D metric.
Same sphere.
\end{proof}

\section{The Propagator Identity}

\begin{theorem}[One propagator, two readings]\label{thm:propagator}
The cascade propagator $K(t) = |K|\,e^{-iHt}$ with $H = \lambda_\infty > 0$
is simultaneously:
\begin{enumerate}
\item the quantum time-evolution operator: unitary, probability-preserving
(Paper~II, Theorem~7.1), and
\item the source of Lorentzian signature: oscillatory character requires
$g_{tt} < 0$ (Paper~III, Theorem~10.2).
\end{enumerate}
\end{theorem}

\begin{proof}
Paper~II derives $K(t)$ from the forced precession $\alpha = \pi/2$ and
identifies $t$ with the slicing direction. Paper~III takes the same $K(t)$
and observes: $e^{-iHt}$ oscillates (unitary), so $g_{tt} < 0$
(Lorentzian). The chain is:
\[
\text{Forced precession} \;\to\; K = |K|\,e^{-iHt}
\;\to\; \text{unitarity (QM)} + \text{Lorentzian (GR)}.
\]
One propagator. One chain. Two names for the same consequence.
\end{proof}

\begin{corollary}[Unitarity and Lorentzian signature are the same theorem]
\label{cor:same}
Both require $H = \lambda_\infty > 0$. This is established once, from the
forced precession, which is established once, from orthogonality. The
cascade's single axiom forces both simultaneously.
\end{corollary}

\begin{remark}[The Euclidean/Lorentzian duality]
The cascade is Euclidean ($B^d$ is a Euclidean object). The physical
spacetime is Lorentzian. The map between them is Wick rotation
$x = it$, but the operation is \emph{cascade-motivated} rather than
\emph{mathematically imposed} (Paper~III~\cite{paper3},
Remark~``Wick rotation, cascade-motivated''): the cascade's forced
precession produces an oscillatory propagator $e^{-i\lambda_\infty x}$ in
real $x$, which fits the Lorentzian regime, and the identification of
$x$ with the observer's proper time requires a Lorentzian sign
$g_{tt} < 0$ on the time-direction metric. In standard QFT, Wick
rotation is a computational trick to render Lorentzian path integrals
convergent; in the cascade, the rotation is the cascade-internal
identification step, motivated by the structure of the propagator.
Both Euclidean and Lorentzian readings are valid: the Euclidean one
is the cascade's starting mathematical structure (the slicing
geometry), the Lorentzian one is the observer's experience (oscillatory
unitary evolution), and the Wick rotation is the bridge.
\end{remark}

\section{The Dirac Descent}\label{sec:dirac-descent}

Theorem~\ref{thm:propagator} establishes that the cascade propagator
$K(t) = |K|\,e^{-iHt}$ is simultaneously the unitary time-evolution
operator of QM and the source of Lorentzian signature. Applied to the
cascade Hilbert space's spinor-representation content
(Paper~III~\cite{paper3}, Lemma~\ref{paper3:lem:cascade-spinor-id}), the
same propagator structure forces a third standard physics result: the
Dirac equation on the chirality bundle in $4$D. The argument runs
along the same template as Theorem~\ref{thm:double}---a standard
mathematical/physical uniqueness result (Wigner's classification of
relativistic wave equations) applied to cascade-derived inputs forces
the standard physics conclusion.

\begin{theorem}[Cascade Dirac descent]\label{thm:cascade-dirac-descent}
At $d=4$, the cascade structure forces the first-order Clifford
equation
\[
(i\gamma^\mu D_\mu - m_f)\,\psi = 0
\]
on the chirality bundle $\mathcal{H} = S^+ \oplus S^-$, where
$D_\mu = \partial_\mu - ieA_\mu$. Each ingredient is cascade-derived:
the gamma matrices $\gamma^\mu$ from Bott classification of the
cascade Clifford algebra $\mathrm{Cl}(1,3)$
(Paper~III~\cite{paper3}, \S\,Clifford classification); the
per-source mass $m_f = R(d_f)/\chi$ from the cascade Berezin
partition with chirality halving at the fermion's Dirac source layer
$d_f$ (Paper~IVb~\cite{paper4b}, Theorem~\ref{paper4b:thm:obstruction});
the charge $e = \sqrt{\alpha(14)}$ from the cascade U(1) gauge
structure at the gauge-window boundary $d_g = 14$
(Paper~IVa~\cite{paper4a}); and the gauge field $A_\mu$ from the
cascade U(1) descent to $4$D.
\end{theorem}

\begin{proof}
Four cascade inputs combine via Wigner-type classification.

\emph{Step~1 (cascade Hilbert space is the spinor representation).}
By Paper~III's Lemma~\ref{paper3:lem:cascade-spinor-id} together with
Theorem~\ref{paper3:thm:d4}, the cascade Hilbert space $\mathcal{H}$ at
the $4$D observer's frame is identified with the irreducibly-complex
minimal spinor representation of $\mathrm{Spin}(1,3)$, equivalently
the $4$-component complex Dirac spinor representation. The chirality
grading $\gamma^5 = i\gamma^0\gamma^1\gamma^2\gamma^3$ supplies the
algebraic decomposition $\mathcal{H} = S^+ \oplus S^-$ into two
$2$-dimensional Weyl basins of equal weight.

\emph{Step~2 (first-order time evolution from the cascade
propagator).} By Theorem~\ref{thm:propagator}, the cascade propagator
$K(t) = |K|\,e^{-iHt}$ generates first-order unitary time
evolution: $i\,\partial_t\psi = H\psi$ for some Hermitian operator
$H$ on $\mathcal{H}$.

\emph{Step~3 (Lorentz covariance forces first-order spatial
derivatives).} A first-order time-derivative equation on a
$\mathrm{Spin}(1,3)$ representation is Lorentz-covariant only if
spatial derivatives are also first-order. Reason: under a boost,
$\partial_t$ mixes with $\partial_i$, so a second-order spatial
structure would, in the boosted frame, produce a second-order time
derivative not present in the original equation, breaking covariance.
By the standard Wigner classification of relativistic wave equations
on a complex spinor representation, the unique Lorentz-covariant
first-order operator is $\gamma^\mu\partial_\mu$ with
$\{\gamma^\mu,\gamma^\nu\} = 2\eta^{\mu\nu}$ (the Dirac matrices,
fixed up to representation choice by Bott classification at
$d \bmod 8 \in \{4,5,6\}$). The equation of motion has covariant
form $(i\gamma^\mu\partial_\mu - m_f)\psi = 0$ for some mass scalar
$m_f$.

\emph{Step~4 (cascade-derived mass and gauge coupling).} By
Paper~IVb's Theorem~\ref{paper4b:thm:obstruction} and the Berezin partition
derivation, the cascade fermion at Dirac source layer $d_f$
($d_f \bmod 8 = 5$) carries mass $m_f = R(d_f)/\chi$, with the mass
operator off-diagonal on the chirality decomposition
$\mathcal{H} = S^+\oplus S^-$ (mass couples chirality basins; this
matches the action of $\beta = \gamma^0$ in the Hamiltonian form).
By Paper~IVa, the cascade U(1) gauge structure at $d_g = 14$
descends to a local U(1) symmetry
$\psi(x)\to e^{i\alpha(x)}\psi(x)$ on $\mathcal{H}$. Local gauge
invariance applied to the first-order operator
$i\gamma^\mu\partial_\mu$ forces the standard minimal-coupling
replacement $\partial_\mu \to D_\mu = \partial_\mu - ieA_\mu$ with
$e = \sqrt{\alpha(14)}$ cascade-derived.

Steps 1--4 give $(i\gamma^\mu D_\mu - m_f)\psi = 0$ at $d=4$ with
all ingredients cascade-derived.
\end{proof}

\begin{lemma}[$4$D-projected fermion field structure]\label{lem:4d-fermion-projection}
Under the cascade structure at $d=4$ (foliation by cross-sections at
height $x \in [-1,1]$ per Paper~III~\cite{paper3}~\S4.1; spinor
representation of $\mathrm{Spin}(1,3)$ on the cascade Hilbert space
per Lemma~\ref{paper3:lem:cascade-spinor-id}), the cascade fermion field
is a section
\[
\psi : M_4 \to \mathcal{S}, \qquad
\mathcal{S} = M_4 \times_{\mathrm{Spin}(1,3)} \mathcal{H},
\]
of the spinor bundle $\mathcal{S}$ over the $4$D Lorentzian manifold
$M_4 = \{ds^2 = -dt^2 + (1-t^2)\,d\sigma_3^2\}$ derived in
Paper~III~\S4.4. The Dirac operator $\gamma^\mu\partial_\mu$ acts
on sections via:
\begin{itemize}[nosep]
\item $\gamma^0\partial_0$: differentiation along the foliation
direction (Paper~III~\S4.1, slicing parameter $x$; Wick-rotated to
the observer's proper time $t$ per Paper~III~\S10).
\item $\gamma^i\partial_i$ ($i = 1,2,3$): differentiation along the
cross-section $S^3$ via the spin-connection lift of the Levi-Civita
connection on $S^3$.
\end{itemize}
The continuous structure on $M_4$ is inherited from the
\emph{continuous} slicing parameter $x$, not from a continuum limit
of the cascade lattice on layer index $d$. The cascade lattice
indexes the tower of dimensions, with $d=4$ being the observer's
layer.
\end{lemma}

\begin{proof}
By Paper~III~\S4.1 the slicing recurrence foliates $B^{d+1}$ into
cross-sections at continuous height $x \in [-1, 1]$, with
$(d+1)$-dimensional metric $ds^2 = dx^2 + (1-x^2)\,d\sigma_{d-1}^2$.
At $d=4$, Paper~III~\S4.4 projects to the FRW metric
$ds^2_{4\mathrm{D}} = -dt^2 + (1-t^2)\,d\sigma_3^2$ via the Wick
rotation $x = it$ of Paper~III~\S10. Call this Lorentzian manifold
$M_4$.

By Lemma~\ref{paper3:lem:cascade-spinor-id}, the cascade Hilbert space
at $d=4$ is the irreducibly-complex minimal spinor representation
$\mathcal{H}$ of $\mathrm{Spin}(1,3)$. The local
$\mathrm{Spin}(1,3)$ action on $M_4$'s tangent frame at each point
lifts to an action on the associated bundle
$\mathcal{S} = M_4 \times_{\mathrm{Spin}(1,3)} \mathcal{H}$. A
cascade fermion configuration is a state in $\mathcal{H}$ varying
smoothly across $M_4$, i.e., a section
$\psi : M_4 \to \mathcal{S}$. Smoothness in the foliation direction
follows from the continuity of the slicing parameter $x$;
smoothness along the cross-section follows from the smooth $S^3$
structure of the spatial slice.

The Levi-Civita connection on $M_4$ lifts uniquely to a spin
connection on $\mathcal{S}$ (standard differential geometry on
$\mathrm{Spin}$-manifolds). The Dirac operator
$\gamma^\mu\nabla_\mu$ on this connection is the unique
Lorentz-covariant first-order differential operator on
$\mathcal{S}$, by the Wigner classification used in Step~3 of
Theorem~\ref{thm:cascade-dirac-descent}. In flat-frame coordinates,
$\gamma^\mu\nabla_\mu = \gamma^\mu\partial_\mu$.

The foliation direction $\gamma^0\partial_0$ is the cascade slicing
direction Wick-rotated to proper time. The spatial directions
$\gamma^i\partial_i$ are the standard $\mathrm{Spin}$-connection on
$S^3$ extended by $\mathrm{Spin}(4) \to \mathrm{Spin}(1,3)$
restriction along the foliation. The cascade lattice on layer
index $d$ does not enter $M_4$'s differential structure; it enters
only as the layer label, with $d=4$ identifying $M_4$ as the
observer's leaf.
\end{proof}

\begin{remark}[$4$D structure derivation: closes presentation gap (c1)]
\label{rem:c1-closed}
Lemma~\ref{lem:4d-fermion-projection} establishes the
spinor-bundle structure on the $4$D Lorentzian manifold parallel
to Paper~III~\S4's metric derivation: the continuous $4$D
structure on which $\gamma^\mu\partial_\mu$ acts is now derived,
not implicit. The relation between the cascade
\emph{tower-action} on the layer index and the $4$D continuous
action on $M_4$ is addressed separately in
Lemma~\ref{lem:within-moment-action-decomposition} and
Remark~\ref{rem:tower-vs-4d-resolved} (genuinely complementary by
design).
\end{remark}

\begin{lemma}[Berezin mass equals Hamiltonian $\beta$-mass]\label{lem:berezin-equals-beta-mass}
The cascade fermion mass term derived from the Berezin partition
(Paper~IVb~\cite{paper4b},
Remark~\ref{paper4b:rem:berezin-partition-derivation}) equals the
Hamiltonian-form $\beta$-mass term in the cascade Dirac equation
of Theorem~\ref{thm:cascade-dirac-descent}. Both act as the
off-diagonal chirality-mixing operator on
$\mathcal{H} = S^+ \oplus S^-$ with strength
$m_f = R(d_f)/\chi$.
\end{lemma}

\begin{proof}
The Berezin fermion action is $S_f = m_f\,\bar\psi\psi$
(Paper~IVb~\cite{paper4b}, Section~\ref{paper4b:sec:fermion-action}, with
$\bar\psi = \psi^\dagger\gamma^0$). On the chirality basis
$\psi = (\psi_+, \psi_-)^T$ in the chiral representation
($\gamma^0$ in $2\times 2$ block form
$\bigl(\begin{smallmatrix} 0 & I \\ I & 0 \end{smallmatrix}\bigr)$,
where $I$ is the $2\times 2$ identity on each Weyl basin), the
scalar pairing reads
\[
\bar\psi\psi \;=\; \psi^\dagger\gamma^0\psi
\;=\; \psi_+^\dagger\psi_- + \psi_-^\dagger\psi_+.
\]
The Berezin mass term $m_f\bar\psi\psi$ thus couples $\psi_+$ to
$\psi_-$ off-diagonally with strength $m_f$.

In Hamiltonian form, the cascade Dirac equation
$(i\gamma^\mu\partial_\mu - m_f)\psi = 0$ rewrites as
$i\partial_t\psi = (\alpha^i p_i + \beta m_f)\psi$, where
$\alpha^i = \gamma^0\gamma^i$ and $\beta = \gamma^0$. The mass
term $\beta m_f = \gamma^0 m_f$ is the same off-diagonal block
matrix in the chiral basis, with strength $m_f$.

Hence the Berezin mass operator $m_f\bar\psi\psi\big|_{\text{op}}$
and the Hamiltonian $\beta$-mass operator $\beta m_f$ are identical:
both equal $\gamma^0\cdot m_f$ acting on
$\mathcal{H} = S^+ \oplus S^-$. The chirality halving
$m = R/\chi$ derived from Poincar\'e--Hopf
(Paper~IVb~\cite{paper4b},
Theorem~\ref{paper4b:thm:obstruction}) and the Hamiltonian-form Dirac
mass agree both in cascade value and in operator action on the
chirality decomposition.
\end{proof}

\begin{remark}[Berezin-mass identification: closes presentation gap (b)]\label{rem:b-closed}
Step~4 of Theorem~\ref{thm:cascade-dirac-descent} parenthetically
identifies the Berezin-mass operator with the Hamiltonian-form
$\beta$-mass. Lemma~\ref{lem:berezin-equals-beta-mass} makes the
identification explicit: both are $\gamma^0\cdot m_f$ on the
chirality basis.
\end{remark}

\begin{remark}[Status of soft spot (a): partial closure with complementary routes]\label{rem:a-closed}
Step~3 of Theorem~\ref{thm:cascade-dirac-descent} invokes Wigner's
classification of relativistic wave equations on a complex spinor
representation to pin $\gamma^\mu\partial_\mu$ as the unique
Lorentz-covariant first-order operator. Soft spot (a) of the
original Open Question asked whether this could be re-derived in
cascade-lattice terms.

Wigner's argument has two distinct contents:
\begin{enumerate}[label=($\alpha$),nosep]
\item \emph{First-order spatial forced by covariance.} Given
first-order time evolution from the cascade propagator
(Theorem~\ref{thm:propagator}, $i\partial_t\psi = H\psi$) and
Lorentz covariance on the spinor representation
(Lemma~\ref{paper3:lem:cascade-spinor-id}), the spatial derivatives
must also be first-order: under a boost, $\partial_t$ mixes with
$\partial_i$, so a second-order spatial operator would, in the
boosted frame, produce second-order time derivatives that the
original first-order equation does not contain.
\end{enumerate}
\begin{enumerate}[label=($\beta$),nosep]
\item \emph{Form of the first-order operator.} Given the
first-order character, the unique Lorentz-covariant first-order
operator on the spinor representation is $\gamma^\mu\partial_\mu$
with $\{\gamma^\mu, \gamma^\nu\} = 2\eta^{\mu\nu}$.
\end{enumerate}

Lemma~\ref{lem:4d-fermion-projection} provides a complementary
\emph{geometric} derivation of part~($\beta$): the spin-connection
lift of the Levi-Civita connection on $M_4$ gives
$\gamma^\mu\nabla_\mu = \gamma^\mu\partial_\mu$ (in flat-frame
coordinates) as the natural first-order operator on the spinor
bundle $\mathcal{S}$, using only cascade-derived inputs.

The spin-connection route does \emph{not} address
part~($\alpha$): the spin connection on a Lorentzian
Spin-manifold is happy with $\gamma^\mu\nabla_\mu$, with the
iterated $(\gamma^\mu\nabla_\mu)^2$, with the scalar Laplacian
$\nabla^\mu\nabla_\mu$ on the spinor bundle, or with any other
covariant operator; the spin-connection structure does not
itself force the equation of motion to be first-order rather
than second-order. Selecting first-order over second-order
requires the cascade-propagator's first-order time form plus
Lorentz covariance on the spinor representation --- exactly
Wigner's content~($\alpha$).

\emph{Status.} Soft spot (a) closes for ($\beta$) by the
spin-connection geometric route in
Lemma~\ref{lem:4d-fermion-projection}. For ($\alpha$), Wigner's
argument remains load-bearing as a standard mathematical theorem
applied to cascade-derived inputs (cascade propagator's
first-order time form per Theorem~\ref{thm:propagator}; Lorentz
covariance on the cascade Hilbert space per
Lemma~\ref{paper3:lem:cascade-spinor-id}). This is the same posture
as Schur's lemma in Paper~III's
Lemma~\ref{paper3:lem:cascade-spinor-id}, Gleason's theorem in
Theorem~\ref{thm:double}, and Lovelock's theorem in
Theorem~\ref{thm:double}: standard mathematical uniqueness
theorems applied to cascade-derived inputs. Wigner's
content~($\alpha$) is therefore at the same rigour level as
those applications, not a substantive defect, but Wigner is
genuinely load-bearing for ($\alpha$) and a stronger ``Wigner
fully superseded'' framing would overstate the closure.
\end{remark}

\begin{lemma}[Within-moment action decomposition]\label{lem:within-moment-action-decomposition}
At fixed observer time $t$, the cascade has two structurally
distinct within-moment geometric directions:
\begin{enumerate}[label=(\roman*),nosep]
\item \emph{Cross-layer (tower depth) direction}: the layer index
$d \in \{4, 5, \ldots, N(t)\}$, indexing the within-moment tower
of dimensions above the observer (Paper~VI~\cite{paper6},
Definition~\ref{paper6:def:tower-time}).
\item \emph{Within-layer ($4$D coordinates) direction}: the
spacetime coordinates $x^\mu$ on $M_4$, the foliated $4$D
Lorentzian manifold at the observer's layer (Paper~III~\S4 +
Lemma~\ref{lem:4d-fermion-projection}).
\end{enumerate}
The cascade's within-moment action splits accordingly:
\begin{itemize}[nosep]
\item \emph{Tower-action}: $S_{\text{tower}}[\varphi] =
\sum_d (2\alpha(d))^{-1}(\Delta\varphi(d))^2$
(Paper~IVb~\cite{paper4b}, Remark~\ref{paper4b:rem:action-uniqueness}),
acting on the cascade tower-scalar $\varphi(d) = \ln\Omega_d$.
\item \emph{$4$D-projected action}: $S_{4\mathrm{D}}[\Phi] =
\int_{M_4} |\partial_\mu\Phi|^2\,d^4x$ (standard Lorentz-covariant
scalar action on $M_4$), acting on $4$D-projected fields
$\Phi(x^\mu)$ derived from cascade-internal content.
\end{itemize}
The two actions act on \emph{different fields} ($\varphi(d)$ on
the tower index vs.\ $\Phi(x^\mu)$ on $M_4$) and parameterize
different cascade-geometric directions (cross-layer vs.\
within-layer). They are therefore genuinely complementary by
design, not two readings of a single underlying action functional.
\end{lemma}

\begin{proof}
The cascade tower-scalar $\varphi(d) = \ln\Omega_d$ is a function
only of the layer index $d$
(Paper~IVb~\cite{paper4b}, Remark~\ref{paper4b:rem:action-uniqueness}).
The tower-action sums layer-difference kinetic
$(\Delta\varphi(d))^2$ across the within-moment tower from
$d = 4$ to $d = N(t)$, normalized by $(2\alpha(d))^{-1}$ from
lattice covariance (Paper~IVb~\cite{paper4b},
Remark~\ref{paper4b:rem:marginal-greens}).

The $4$D-projected action acts on fields
$\Phi : M_4 \to \mathbb{R}$ (or higher-spin analogs;
Lemma~\ref{lem:4d-fermion-projection} for the spinor case). The
form $\int |\partial_\mu\Phi|^2\,d^4x$ is the unique
Lorentz-covariant second-order kinetic action on $M_4$ (standard
field theory on a Lorentzian manifold; same posture as
Theorem~\ref{thm:double}'s use of Lovelock for the gravitational
sector and Theorem~\ref{thm:cascade-dirac-descent}'s use of
Wigner for the spinor sector).

The two action functionals act on functions with disjoint
domains: $\varphi(d)$ has domain $\mathbb{Z}_{\geq 4}$;
$\Phi(x^\mu)$ has domain $M_4$. No single field carries both,
and no single action functional unifies them: the two actions
describe orthogonal cascade directions of the within-moment
geometry.
\end{proof}

\begin{remark}[Tower-action vs.\ $4$D action: genuinely complementary by design]
\label{rem:tower-vs-4d-resolved}
A natural question on the cascade Dirac descent is whether the
cascade tower-action $\sum_d (2\alpha(d))^{-1}(\Delta\varphi(d))^2$
and the standard $4$D continuous action
$\int|\partial\varphi|^2$ are genuinely complementary (a), unifiable
via a foliated total-space action (b), or unifiable via the Wick
rotation reading $\Delta\varphi(d) \leftrightarrow \partial_t\Phi$ (c).
Lemma~\ref{lem:within-moment-action-decomposition} establishes
reading~(a) as the resolution, with structural justification.

The two actions act on \emph{different fields} (tower-scalar
$\varphi(d)$ vs.\ $4$D-projected $\Phi(x^\mu)$) and govern
\emph{different cascade-geometric directions} (cross-layer vs.\
within-layer). Reading~(b) is structurally incorrect: there is no
single field on which both actions act, hence no functional that
sums to give both as projections. Reading~(c) is structurally
incorrect: the layer index $d$ is the cascade's tower-depth
parameter, not a Wick-rotation of $4$D time --- $4$D time arises
from Wick-rotating the within-layer foliation parameter
$x \in [-1,1]$ (Paper~III~\S10, $x = it$), which is structurally
distinct from the layer index. The cascade's two distinct
time-like / spatial-like cascade directions
(Paper~VI~\cite{paper6}, \S1.2 within-moment vs
between-moments) provide the structural framework: the
within-moment tower-action governs cross-layer kinetics
(direction~(i) of
Lemma~\ref{lem:within-moment-action-decomposition}); the
within-moment $4$D action governs within-layer kinetics
(direction~(ii)); the between-moments tower-growth operator
$A$ (Paper~VI~\cite{paper6},
Proposition~\ref{paper6:prop:time-evolution}) governs the genuine
cosmic-time dynamics by adding layers at the top of the tower.

This is the cascade's complete within-moment + between-moments
action architecture. Each piece governs a structurally distinct
cascade direction; none unifies into a single functional, but
each is independently derivable from cascade primitives. The
tower-action is derived in Paper~IVb~\cite{paper4b}'s
Remark~\ref{paper4b:rem:action-uniqueness} from lattice covariance;
the $4$D-projected action follows from standard Lorentz
covariance on the cascade-derived $M_4$
(Paper~III~\S4); the between-moments tower-growth $A$ is
defined in Paper~VI~\cite{paper6}.
\end{remark}

\begin{remark}[Status: parallel to Theorem~\ref{thm:double}]
\label{rem:cascade-dirac-status}
Theorem~\ref{thm:cascade-dirac-descent} sits at the same rigour
level as Theorem~\ref{thm:double}: a standard mathematical
uniqueness result (Wigner's classification, replacing
Gleason$+$Lovelock) applied to cascade-derived inputs (cascade
Hilbert space, chirality decomposition, mass, gauge structure)
forces the standard physics result (the Dirac equation, replacing
QM/GR consistency). The cascade contribution is to provide the
inputs; the standard mathematical theorem is invoked, not
re-derived---the same posture Paper~III adopts when invoking
Schur~$+$~Clifford classification in
Lemma~\ref{paper3:lem:cascade-spinor-id}, and the same posture
Theorem~\ref{thm:double} adopts when invoking Gleason and Lovelock.
Presentation gaps in Steps~3 and~4 are addressed by
Lemma~\ref{lem:4d-fermion-projection},
Lemma~\ref{lem:berezin-equals-beta-mass}, and
Remarks~\ref{rem:c1-closed},~\ref{rem:b-closed},~\ref{rem:a-closed}
above.
\end{remark}

\begin{remark}[Triple uniqueness from a single source]
\label{rem:triple-uniqueness}
The series now exhibits three instances of the same template
``cascade single source $+$ standard uniqueness theorem
$\Rightarrow$ standard physics result'':
\begin{itemize}[nosep]
\item \emph{Born rule of QM}: cascade Hilbert space $+$ Gleason
$\Rightarrow$ probability measure unique
(Theorem~\ref{thm:double}, clause~2).
\item \emph{Einstein equation of GR}: cascade $4$D Lorentzian
manifold $+$ Lovelock $\Rightarrow$ gravitational equation unique
(Theorem~\ref{thm:double}, clause~3).
\item \emph{Dirac equation of matter dynamics}: cascade
spinor-rep Hilbert space $+$ Wigner classification $\Rightarrow$
relativistic spin-$\tfrac{1}{2}$ wave equation unique
(Theorem~\ref{thm:cascade-dirac-descent}).
\end{itemize}
The cascade's claim that ``the infinite-dimensional unit ball,
descended to $4$D, is indistinguishable from our universe''
acquires an additional confirmation: the $4$D matter dynamics,
like the QM and gravitational projections, is forced once the
cascade single source is fixed. The double-uniqueness picture of
Theorem~\ref{thm:double} extends to a triple-uniqueness picture
covering QM dynamics, gravitational dynamics, and matter dynamics
under one structural template.
\end{remark}

\section{Double Uniqueness}

\begin{theorem}[Gleason $+$ Lovelock $+$ single source $=$ consistency]\label{thm:double}
At $d = 4$, the cascade's quantum and gravitational predictions on any
shared 4D observable agree. The argument combines three inputs:
\begin{enumerate}
\item Single-source structure (Theorem~\ref{thm:single}): both
projections $\mathcal{Q}$ and $\mathcal{G}$ are deterministic
functions of the same cascade state $|\Psi\rangle\in S^{d-1}$.
\item Gleason's theorem~\cite{gleason}: given the cascade-derived
Hilbert space (Paper~II), the Born rule is the unique probability
measure on its lattice of subspaces (for Hilbert dimension $\geq 3$).
\item Lovelock's theorem~\cite{lovelock}: given the cascade-derived
4D Lorentzian manifold (Paper~III), the Einstein equation is the unique
symmetric divergence-free rank-2 tensor equation in $d = 4$.
\end{enumerate}
\end{theorem}

\begin{proof}
Let $\mathcal{O}$ be any observable computable from both the quantum
projection $\mathcal{Q}$ and the gravitational projection
$\mathcal{G}$ of a cascade state $|\Psi\rangle$. By
Theorem~\ref{thm:single}, $\mathcal{Q}(|\Psi\rangle)$ and
$\mathcal{G}(|\Psi\rangle)$ are both determined by $|\Psi\rangle$ with no
additional structure.

By Gleason's uniqueness, among all probability measures on the
cascade-derived Hilbert space's subspace lattice, only the Born rule is
admissible; therefore the quantum prediction for $\mathcal{O}$,
$\mathcal{Q}(|\Psi\rangle)(\mathcal{O})$, is not a choice but a
computation from $|\Psi\rangle$. By Lovelock's uniqueness, among all
gravitational equations on the cascade-derived 4D manifold, only the
Einstein equation is admissible; therefore the gravitational prediction
for $\mathcal{O}$, $\mathcal{G}(|\Psi\rangle)(\mathcal{O})$, is also a
computation from $|\Psi\rangle$.

Consistency $\mathcal{Q}(|\Psi\rangle)(\mathcal{O})
= \mathcal{G}(|\Psi\rangle)(\mathcal{O})$ on any shared observable
follows because both sides are computations of the same cascade
quantity---an observable of $|\Psi\rangle$---via two fixed, unique
projections with no freedom to produce differing answers from the same
input. The two projections compute the same thing through different
routes; Gleason and Lovelock ensure neither route admits a variant that
would predict a different value.
\end{proof}

\begin{remark}[What the argument does \emph{not} claim]\label{rem:sp26-status}
The argument is \emph{not} that ``each theory is unique in its own
domain, therefore they cannot disagree on shared observables.'' That
inference is invalid: two theories each individually unique in
separate domains can disagree on shared observables. Quantum
electrodynamics is unique given gauge invariance; effective-field-theory
quantum gravity is unique given the EFT postulates; yet their
predictions on quantum-gravitational scattering amplitudes diverge at
high energies because the two theories do not share a common source and
their domains overlap only by heuristic matching.

The cascade's consistency claim is stronger and rests on
Theorem~\ref{thm:single}: the quantum and gravitational projections are
not two independent theories that happen to be unique in their own
domains, they are two readings of the same cascade geometry. Shared
observables are computed from one input ($|\Psi\rangle$) via
deterministic projections; Gleason and Lovelock rule out alternative
projections with different outputs, and the single source
(Theorem~\ref{thm:single}) rules out the QED/EFT-QG scenario---there
is no ``separately derived'' alternative theory on either side of the
cascade. Without Theorem~\ref{thm:single}, the two uniqueness theorems
alone would not imply cross-domain consistency. This is the substantive
content of the ``double uniqueness'' label: Gleason and Lovelock lock
each projection in place, and the single-source architecture guarantees
both projections read the same source.
\end{remark}

\begin{remark}[The role of $d = 4$]
The argument depends on $d = 4$. In higher dimensions, Lovelock permits
additional gravitational terms with free couplings; the gravitational
projection would no longer be unique and consistency would fail. The
cascade forces $d = 4$ (Paper~III, Section~9) precisely where both
uniqueness theorems hold.
\end{remark}

\section{No Absolute Scale}

\begin{theorem}[The cascade produces ratios]\label{thm:noscale}
Every physical prediction of the cascade series is a dimensionless ratio.
The cascade geometry contains no intrinsic length, time, mass, or energy
scale. The dimensionful constants $G$, $\hbar$, $c$ are unit-conversion
factors, not physical parameters.
\end{theorem}

\begin{proof}
The cascade's input is $B^d$ with $|x| \leq 1$. Rescaling changes no
ratio. Its outputs---$I$, $\alpha_s$, $\sin^2\theta_W$, $m_\mu/m_e$,
$\Omega_m$---are pure numbers from the Gamma function. $\hbar = (\pi/2)/\Delta t$ (Paper~II): the angle is derived; $\Delta t$ is a unit
choice. $G$ maps cascade curvature ratios to human units. $c$ converts
spatial and temporal components of the cascade metric.
\end{proof}

\begin{remark}[Cascade-intrinsic vs.\ empirically-anchored scales]\label{rem:sp24-status}
Theorem~\ref{thm:noscale} makes a statement about the cascade geometry
in isolation: it produces dimensionless ratios and no intrinsic
scales. This is compatible with Part~IVb's absolute-mass predictions
($m_\tau = 1777$~MeV, $m_W = 80.10$~GeV, $v = 240.8$~GeV, etc.) once
the distinction between \emph{cascade-intrinsic} and
\emph{empirically-anchored} quantities is made explicit:
\begin{itemize}
\item \emph{Cascade-intrinsic (dimensionless, zero empirical input).}
All ratios listed above, and every other prediction of the series
that enters as a ratio: $\alpha_s$, $\sin^2\theta_W$, $m_\mu/m_e$,
$\Omega_m$, $1/\alpha_{\mathrm{GUT}}$, $\theta_C$, and so on. These
are Gamma-function values; no external scale enters.
\item \emph{Empirically-anchored (dimensional).} Absolute quantities
with units (MeV, GeV, Mpc, K, $\ldots$). Each requires exactly one
empirical input per independent sector: Part~IVb takes
$M_{\mathrm{Pl,red}} = 2.435\times 10^{18}$~GeV as the gravitational
anchor and $m_Z = 91.19$~GeV as the electroweak anchor
(Remark after Theorem~4.4 in Part~IVb). All other absolute masses
follow as cascade-derived ratios multiplied by these anchors; no
further empirical input is required.
\end{itemize}
This is the standard position of physical theories: general
relativity predicts the form of the Schwarzschild metric and uses
Newton's $G$ as the gravitational anchor; quantum electrodynamics
predicts scattering cross-sections as functions of $\alpha$ but takes
$\alpha$ itself from measurement. The cascade's claim is not ``no
dimensional numbers appear in absolute predictions,'' which would be
absurd, but ``no dimensional number is intrinsic to the cascade
geometry.'' Plugging $M_{\mathrm{Pl,red}}$ and $m_Z$ into the
cascade's dimensionless ratios to recover the full dimensional mass
spectrum is a valid sanity check, not a circularity: the cascade
produces the ratios, the anchors supply the units, and the resulting
agreement with observation (e.g., $m_\tau = 1776.82$~MeV derived from
$M_{\mathrm{Pl,red}}$ via Theorem~4.10 of Part~IVb) is a testable
prediction rather than a fit---the only freedom is the anchor value,
which is fixed by a single external measurement per sector. Compare
with the Standard Model, which has $\sim 22$ independent dimensional
parameters; the cascade has $2$.
\end{remark}

\begin{corollary}[The UV problem is a projection artefact]
The ``Planck scale'' is not a physical place. It is a ratio at which the 4D
effective description becomes unreliable. The cascade remains well-defined
at all ratios. Standard UV divergences arise from integrating over
arbitrarily high momenta in 4D; in the cascade, the sum over levels
terminates at $d = 217$. The finiteness is structural (a theorem about the
Gamma function), not regulatory.
\end{corollary}

\section{Metric Superpositions Are Cascade States}

\begin{theorem}[Gravity is already quantum]\label{thm:metricstate}
The 4D metric $g_{\mu\nu}$ and the cascade state
$|\Psi\rangle \in S^{d-1}$ are both cascade-geometric quantities: the
metric is derived from the cascade's foliation (Paper~III, \S4), and the
state is a point on the cascade's boundary sphere (Paper~II,~\S2). The
cascade's single source therefore predicts that metric structure and
quantum state are coupled: different cascade states should produce
different metrics, and superpositions of states should produce
superpositions of metrics. Under this coupling, the Born rule (Gleason)
assigns probabilities to metric outcomes and each outcome satisfies
Einstein's equation (Lovelock).
\end{theorem}

\begin{proof}
Three observations, each established elsewhere in the series.

\emph{Step~1 (derived).} Paper~III~\S4 derives $g_{\mu\nu}$ from the
cascade's slicing recurrence as an FRW-type metric
$ds^2 = -dt^2 + a(t)^2\,d\sigma_3^2$ with $a(t) = \sqrt{1-t^2}$
(Lorentzian cross-section radius, Paper~III Theorem~9.1).
The metric is a cascade-geometric object.

\emph{Step~2 (derived).} Paper~II~\S2 establishes the cascade state space
as $S^{d-1}$ with superposition structure $|\Psi\rangle = a|\Psi_1\rangle
+ b|\Psi_2\rangle$. States are cascade-geometric objects.

\emph{Step~3 (asserted, pending a concrete state-metric construction).}
That different states produce distinct metrics---i.e., a bijection or
more generally a well-defined map
$|\Psi\rangle \mapsto g_{\mu\nu}^{(\Psi)}$ with
$g^{(1)}\neq g^{(2)}$ for at least one pair
$|\Psi_1\rangle,|\Psi_2\rangle$---is asserted on single-source grounds
(both objects derive from the cascade; their joint theory is not two
independent sectors) but has not been instantiated in Paper~III. The FRW
metric of \S4 is a single background; the paper does not exhibit a pair
of cascade states mapping to distinct metrics. The Born rule on metric
outcomes and the Lovelock uniqueness of each outcome follow
conditionally on Step~3's concrete construction. See
Remark~\ref{rem:sp25-status} for derivation status and
Open Question~\ref{oq:state-metric-instantiation} for the cascade-native
construction target.
\end{proof}

\begin{remark}[Derivation status of metric superpositions]\label{rem:sp25-status}
The theorem's three claims factorise with three different derivation
statuses:
\begin{enumerate}
\item \emph{Derived.} The metric $g_{\mu\nu}$ is a cascade-geometric
quantity (Paper~III~\S4, slicing foliation).
\item \emph{Derived.} The cascade state $|\Psi\rangle$ is a
cascade-geometric quantity (Paper~II~\S2, point on $S^{d-1}$).
\item \emph{Asserted, pending construction.} The map
$|\Psi\rangle \mapsto g_{\mu\nu}^{(\Psi)}$ with state-dependent
output---the core content of ``different states produce different
metrics''---is stated on single-source grounds but not exhibited
concretely. Paper~III's derivation yields a single FRW background
$-dt^2 + (1-t^2)\,d\sigma_3^2$ that is the same for all cascade states
on $S^{d-1}$. Two plausible cascade-native routes to a concrete
construction are: (a) state-dependent foliation---different states
select different time directions in $B^{d+1}$'s embedding, producing
distinct Lorentzian slicings of the same cascade geometry; and (b)
stress-energy back-reaction---different states carry different
cascade-content distributions, sourcing different matter terms on the
right-hand side of Einstein's equation via the projection identity of
Paper~III~\S8 and producing distinct metrics as solutions. Neither
route is worked out in Paper~III. The paper's ``Does not'' list acknowledges the
related but adjacent gap (``construct the explicit Hilbert space of
metric superpositions or compute decoherence rates'').
\end{enumerate}
Under the conditional in Step~3, Gleason's Born rule gives probabilities
to metric outcomes and Lovelock gives Einstein's equation per outcome;
these steps are the content of Theorem~\ref{thm:double}.
Closing Step~3 would complete the ``gravity is already quantum'' claim
by exhibiting a concrete cascade-native map from states to metrics and
thereby concretising the metric-superposition interpretation.
\end{remark}

\begin{remark}[No second quantisation of gravity]
The standard approach: take $g_{\mu\nu}$, promote to $\hat{g}_{\mu\nu}$,
impose commutation relations, struggle with non-renormalisability. The
cascade bypasses this. Under the coupling of Theorem~\ref{thm:metricstate}
(pending the Step~3 construction of
Remark~\ref{rem:sp25-status}), the metric is never promoted to an
operator: it is a property of the quantum state. Superpositions of
metrics are superpositions of states. Gravity was always quantum in the
cascade, because geometry was always quantum in the cascade. The
dissolution is structural; the concrete state-metric map remains to be
written.
\end{remark}

\section{Bekenstein--Hawking Entropy Is Hidden Geometry}

\subsection{The claim}

A black hole's ``entropy'' is not a statistical quantity. It is the amount of
cascade geometry the observer cannot see.

\begin{theorem}[Hidden geometry behind a horizon]\label{thm:bh}
Let a $d$-dimensional cascade observer measure a horizon of boundary area $A$
in cascade units. The geometric content hidden behind the horizon is
\[
\boxed{\;S = \frac{A}{d}.\;}
\]
At $d = 4$: $S = A/4$.
\end{theorem}

\begin{proof}
Three steps.

\emph{Step 1 (derived, unit ball).} Boundary dominance
(Paper~0, Theorem~3.1~\cite{paper0}, \ref{paper0:thm:boundary}):
$\Omega_{d-1}/V_d = d$, hence $V_d = \Omega_{d-1}/d$.

\emph{Step 2 (derived, cascade-content primitive).} By Paper~0,
Corollary~3.2 (\ref{paper0:cor:primary}), sphere area is the unique independent
cascade quantity: content~$\equiv$~area, and every volume is a derived
quantity. The unit-ball identity therefore reads
``content inside a unit sphere of area $\Omega_{d-1}$ is $\Omega_{d-1}/d$.''

\emph{Step 3 (asserted, scale invariance of the content--area map).}
For a cascade horizon of area $A$ in cascade units, we \emph{assert} the
map from area to hidden content is linear in $A$ with slope $1/d$, so that
$S = A/d$ for any $A$, not only $A = \Omega_{d-1}$. Taken as a
geometric radius-Euclidean identity this step does not follow from Step~1
alone: for a ball of physical radius $r$ at layer $d$, $V(r)/A(r) = r/d$,
so $1/d$ is recovered only at $r=1$. Here we are instead claiming that the
cascade counts content by area-units (each unit of cascade area contributing
$1/d$ units of hidden content) and that this counting is additive and
scale-free. See Remark~\ref{rem:sp23-status} for the derivation status of
this step and §\ref{sec:open-questions-part2eq3},
Question~\ref{oq:content-area-scale-invariance}, for the cascade-native
derivation target. The empirical confirmation at $d=3$ (BTZ,
Theorem~\ref{thm:btz}) and at $d=4$ (Schwarzschild, Theorem~\ref{thm:hawking})
closes the formula on observation; the derivation gap is mechanistic, not
numerical.

At $d = 4$: $S = A/4$.
\end{proof}

One theorem. Three steps. No semiclassical gravity, no QFT on curved spacetime,
no thermal radiation, no microstate counting.

\begin{remark}[Derivation status of $S = A/d$]\label{rem:sp23-status}
The entropy formula factorises into three ingredients with three different
derivation statuses:
\begin{enumerate}
\item \emph{Derived.} The unit-ball boundary-dominance identity
$V_d = \Omega_{d-1}/d$ (Paper~0, Theorem~3.1~\cite{paper0}). A Gamma-function
identity; no assumption.
\item \emph{Derived as a cascade primitive.} Content~$\equiv$~area and
volumes are derived: Paper~0, Corollary~3.2 (\ref{paper0:cor:primary}).
No assumption beyond the cascade's definition of independent quantities.
\item \emph{Asserted, pending a cascade-action derivation.} That the
content--area map is linear in $A$ with slope $1/d$ for \emph{any} cascade
area, not only for $A = \Omega_{d-1}$. Radius-Euclidean geometry alone gives
$V(r)/A(r) = r/d$ for a ball of physical radius $r$, so the cascade's
scale-free identity $S/A = 1/d$ requires an additional counting principle:
cascade content is accumulated additively in units of $V_d/\Omega_{d-1}$ per
unit area. This principle is the operational content of
Corollary~3.2's ``sphere area is the unique independent cascade quantity''
taken as a quantitative identification rather than a qualitative one.
\end{enumerate}
Empirical support for the asserted Step~3 is strong: the same formula
$S = A/d$ reproduces Bekenstein--Hawking at $d=4$ (Theorem~\ref{thm:hawking})
and BTZ thermodynamics at $d=3$ (Theorem~\ref{thm:btz}, with
$G_3 = 3/4$ from Theorem~\ref{thm:Gd}, verified numerically in
\texttt{tools/closures/btz\_cross\_check.py}) without parameter fitting. The
residual is mechanistic: the numerical relation is correct at two
independent layers; what is missing is a cascade-action derivation that
forces the linear content--area map without invoking the additivity
principle as input.
\end{remark}

\subsection{Content is geometry}

Corollary~3.2 of Paper~0 (\ref{paper0:cor:primary}): sphere area is the only
independent cascade quantity. Content IS geometry. One unit of sphere area
is one unit of cascade content. A horizon of area $A$ hides $A/4$ units of
geometric content. The question ``how many microstates?'' is
downstream---the geometry is primary.

\subsection{Why Hawking got \texorpdfstring{$A/4$}{A/4} from a different calculation}

Hawking (1975)~\cite{hawking75} derived $S = A/4$ from QFT on a Schwarzschild
background: Bogoliubov transformations, thermal density matrices, the
semiclassical approximation. A statistical calculation producing a geometric
answer. This has been mysterious for fifty years.

The cascade explains: Hawking measured the hidden geometry through a 4D
effective theory. The thermal radiation is the 4D observer's description of
geometric content leaking across an asymptotic compactification boundary.
Both routes arrive at $A/4$ because they measure the same geometry from
different altitudes.

\subsection{The standard interpretations are downstream}

\begin{center}
\begin{tabular}{ll}
\hline
Level & Statement \\
\hline
Cascade (primary) & Hidden geometry $= A/d$ \\
$\downarrow$ & \\
Thermodynamic & $Q/T = A/4$ \\
Statistical & $\ln W = A/4$ \\
Information-theoretic & $H = A/4$ nats \\
\hline
\end{tabular}
\end{center}

All three give $A/4$ because they are all measuring the same geometric
content. None is fundamental. The geometry is fundamental.

\subsection{Hawking temperature from the cascade}

\begin{theorem}[Hawking temperature]\label{thm:hawking}
A Schwarzschild black hole of mass $M$ has temperature
\[
T = \frac{1}{8\pi M}.
\]
\end{theorem}

\begin{proof}
Three steps, each from the cascade's own content.

\emph{Step~1.} Hidden geometry $S = A/4$ (Theorem~\ref{thm:bh}, boundary
dominance at $d = 4$).

\emph{Step~2.} The Einstein equation is the unique gravitational equation at
$d = 4$ (Paper~III, Lovelock uniqueness). By Birkhoff's theorem---a
mathematical consequence of the Einstein equation---the Schwarzschild metric
is the unique spherically symmetric vacuum solution. Its horizon has area
$A = 16\pi M^2$ in Planck units.

\emph{Step~3.}
\[
T = \left(\frac{\partial S}{\partial M}\right)^{-1}
= \left(\frac{\partial}{\partial M}\frac{A}{4}\right)^{-1}
= \left(\frac{\partial}{\partial M}\,4\pi M^2\right)^{-1}
= \frac{1}{8\pi M}.\qedhere
\]
\end{proof}

\begin{remark}[What the derivation uses and does not use]
The derivation uses boundary dominance (a Gamma function identity),
Lovelock uniqueness (a tensor-algebraic theorem), Birkhoff's theorem
(a mathematical consequence of the Einstein equation), and the chain rule.
It does not use quantum field theory on curved spacetime, Bogoliubov
transformations, the semiclassical approximation, or any counting of
microstates. Hawking's original derivation~\cite{hawking75} required
semiclassical gravity because it had no geometric route to $S = A/4$; the
cascade provides that route through boundary dominance, making the
semiclassical calculation unnecessary.
\end{remark}

\begin{remark}[The first law is output, not input]
The relation $dM = T\,dS$ is often cited as the first law of black hole
mechanics and used as an input to derive $T$. Here the logic is reversed:
$S(M) = 4\pi M^2$ is known from Steps~1--2, so $T = (dS/dM)^{-1}$ follows
by calculus. The first law is the \emph{output}---a consequence of the
cascade's geometric entropy and the Einstein equation---not an independent
postulate.
\end{remark}

\subsection{Cosmological hidden geometry}

The cascade hierarchy $\Omega_7/\Omega_{217} \approx 10^{121}$ is the ratio
of sphere areas at the two equilibria: the largest and the smallest in the
cascade's tower. At $d = 4$, boundary dominance gives the hidden geometry
as $1/4$ of this hierarchy---approximately $10^{120}$ in cascade units.
This is the same $10^{120}$ that appears as the cosmological constant's
hierarchy (Paper~I, Theorem~9.1), seen from the other side: $\Lambda \sim
10^{-120}$ is the inverse of the hidden geometry. The cascade derives both
numbers from the Gamma function; neither requires a de~Sitter solution or
semiclassical gravity.

\subsection{BTZ cross-check at \texorpdfstring{$d=3$}{d=3}}
\label{subsec:btz}

Theorem~\ref{thm:bh} applies at any $d$ for which boundary dominance holds,
which is any $d\geq 1$ (Paper~0, Theorem~3.1). It is not restricted to the
observer's dimension. Applied at $d=3$ it predicts
$S = A/3$ for any $d=3$ horizon, and in particular for the BTZ black
hole~\cite{btz}, the unique rotating/stationary solution of $2+1$
Einstein gravity with negative cosmological constant. This provides an
independent structural test of boundary dominance at a dimension other
than the observer's.

\begin{theorem}[Cascade prediction for $G_d$]\label{thm:Gd}
Matching the cascade's boundary-dominance entropy $S = A/d$
(Theorem~\ref{thm:bh}) to standard $d$-dimensional Bekenstein--Hawking
entropy $S_{\mathrm{BH}} = A/(4 G_d)$ requires
\[
G_d \;=\; \frac{d}{4} \quad\text{(in cascade units where $\ell_{\mathrm{Pl}}=1$).}
\]
At $d=4$ this reproduces the standard normalisation $G_4 = 1$ used in
Theorem~\ref{thm:hawking}. At $d=3$ it predicts $G_3 = 3/4$, and therefore
the dimensionless ratio $G_3/G_4 = 3/4$.
\end{theorem}

\begin{proof}
The cascade's formula $S = A/d$ is a geometric identity
(boundary dominance). Identifying this with standard Bekenstein--Hawking
entropy $S_{\mathrm{BH}} = A/(4G_d)$ gives $1/d = 1/(4G_d)$, hence
$G_d = d/4$.
\end{proof}

\begin{theorem}[BTZ thermodynamics from the cascade]\label{thm:btz}
For the non-rotating uncharged BTZ black hole in $2+1$ gravity with AdS
radius $\ell$, with $G_3 = 3/4$ (Theorem~\ref{thm:Gd}):
\begin{itemize}[nosep]
\item horizon: $r_h^2 = 8 G_3 M \ell^2 = 6 M \ell^2$;
\item area (circumference of $S^1$ horizon): $A = 2\pi r_h = 2\pi\ell\sqrt{6M}$;
\item cascade entropy: $S = A/3 = (2\pi\ell/3)\sqrt{6M}$;
\item Hawking temperature from the first law $T = (dS/dM)^{-1}$:
\[
T \;=\; \frac{r_h}{2\pi\ell^2} \;=\; \frac{\sqrt{6M}}{2\pi\ell}.
\]
\end{itemize}
Both $S$ and $T$ coincide exactly with the known BTZ
results~\cite{btz,carlip-btz}: $S_{\mathrm{BTZ}} = A/(4G_3)$ and
$T_{\mathrm{BTZ}} = r_h/(2\pi\ell^2)$ from surface gravity. The cascade's
boundary-dominance entropy formula, applied at $d=3$, is self-consistent
with standard BTZ thermodynamics to all order.
\end{theorem}

\begin{proof}
Direct computation; mechanically verified in the script
\texttt{tools/\allowbreak closures/\allowbreak btz\_cross\_check.py}.
The substitution
$G_3 = 3/4$ converts $S_{\mathrm{BH}} = A/(4G_3)$ into $A/3$ identically.
The first-law temperature
\[
T \;=\; \left(\frac{dS}{dM}\right)^{-1}
\;=\; \left(\frac{d}{dM}\frac{2\pi\ell\sqrt{6M}}{3}\right)^{-1}
\;=\; \left(\frac{\pi\ell\sqrt{6}}{3\sqrt{M}}\right)^{-1}
\;=\; \frac{\sqrt{6M}}{2\pi\ell}
\;=\; \frac{r_h}{2\pi\ell^2},
\]
which is the standard BTZ surface-gravity temperature. No approximation
is made; every step is exact.
\end{proof}

\begin{remark}[What this cross-check establishes]
The BTZ cross-check is a \emph{structural} test of boundary dominance at a
dimension other than the observer's. The Gamma-function identity
$\Omega_{d-1}/V_d = d$ holds for all $d \geq 1$; applying it at $d=3$
produces a complete thermodynamic description (entropy, temperature,
first law) that agrees with standard BTZ without any free parameter. The
non-trivial output is the cascade's prediction $G_d = d/4$, giving
specifically $G_3/G_4 = 3/4$ for the ratio of Newton constants in three
versus four dimensions.

Theorems~\ref{thm:Gd} and~\ref{thm:btz} together constitute a
cascade-internal derivation: $G_d = d/4$ follows from boundary dominance
plus identification with the standard $A/(4G_d)$ entropy formula; both
are cascade-internal quantities (boundary dominance is a Gamma-function
theorem; the $A/(4G_d)$ form is Bekenstein--Hawking applied at dimension
$d$, itself a consequence of the cascade's Einstein equation at that $d$
via Lovelock). No semiclassical integration enters. The cascade's own
structural commitments (Paper~I~\cite{paper1}~\S3.2;
Paper~II\,=\,III~\S5--6; Paper~III~\cite{paper3}~\S12) explicitly
\emph{exclude} Kaluza--Klein reduction as the route to $G_d$: the metric
is a state property, not an operator on a fixed background, so there is
no semiclassical reduction over compactified dimensions to perform.

What would strengthen the cross-check is a \emph{second} cascade-internal
route to $G_d$ that does not go through the matching
$S = A/d \leftrightarrow A/(4G_d)$---for instance, via the proposed
discrete cascade action (Paper~IVb~\cite{paper4b}, Remark~4.6) evaluated
at layer $d$, or via an explicit cascade-native sphere-area identity at
dimension $d$ that independently equals $d/4$. Any such route would
promote Theorem~\ref{thm:Gd} from ``structural matching'' to
``doubly-derived.'' Together with Theorem~\ref{thm:hawking}
(Schwarzschild at $d=4$), Theorem~\ref{thm:btz} is the second confirmed
dimension of the cascade's entropy formula.
\end{remark}

\section{Resolution of the Standard Conflicts}

\begin{center}
\small
\begin{tabular}{p{3.5cm}p{3.8cm}p{4.5cm}}
\hline
Conflict & Source & Cascade resolution \\
\hline
UV catastrophe & Infinite modes in 4D & Finite tower: 213 levels, structural \\
Problem of time & $H = 0$ in canonical gravity & $H = \lambda_\infty > 0$; time is
slicing \\
Background dependence & QM needs fixed metric & Hilbert space on $S^{d-1}$; metric
is output \\
Information paradox & Sharp horizon $+$ unitarity & No sharp horizon; asymptotic
compactification \\
Non-renormalisability & $g_{\mu\nu} \to \hat{g}_{\mu\nu}$ & Metric is state property;
no promotion \\
Planck-scale physics & What happens at $l_{\mathrm{Pl}}$? & No absolute scale; cascade
defined at all ratios \\
Euclidean/\allowbreak Lorentzian & Wick rotation ambiguity & Wick rotation cascade-motivated:
forced precession $\Rightarrow$ oscillatory propagator $\Rightarrow$ $g_{tt} < 0$ \\
BH entropy & Statistical $\to$ geometric? & Geometry is primary: $S = A/d$ \\
\hline
\end{tabular}
\end{center}

Every row uses the cascade's existing content (Papers~I--III) and classical
theorems (Gleason, Lovelock). No new postulate, no new parameter, no new
mathematical structure.

\section{Summary}

\begin{center}
\small
\begin{tabular}{llc}
\hline
Result & Mechanism & Section \\
\hline
QM and GR share single source & Same state space, propagator, hypothesis & 2 \\
Hilbert space $=$ configuration space & Both are $S^{d-1}$ & 2.2 \\
Unitarity $=$ Lorentzian signature & Same propagator, same $H > 0$ & 3 \\
Mutual consistency & Gleason $+$ Lovelock: double uniqueness & 4 \\
No absolute scale & All predictions are ratios & 5 \\
UV dissolution & Finite tower; projection artefact & 5 \\
Gravity already quantum & Metric is state property on $S^{d-1}$ & 6 \\
$S = A/4$ & Boundary dominance at $d = 4$ & 7 \\
$T = 1/(8\pi M)$ & $S(M) + $ Birkhoff; no semiclassics & 7.4 \\
Hidden geometry $\sim 10^{120}$ & Cascade hierarchy $/\, 4$ & 7.5 \\
\hline
\end{tabular}
\end{center}

\section{What This Paper Does and Does Not Do}

\textbf{Does:}
\begin{itemize}
\item Proves the quantum and gravitational projections are the same theorem:
same source, same propagator, same state space, same hypothesis.
\item Proves mutual consistency via double uniqueness (Gleason $+$ Lovelock).
\item Dissolves all seven standard QM/GR conflicts.
\item Derives $S = A/4$ from one theorem (boundary dominance), with no
semiclassical input.
\item Identifies the factor $1/4$ as $1/d$ at the observer's dimension.
\item Derives the Hawking temperature $T = 1/(8\pi M)$ from boundary
dominance and Birkhoff's theorem, with no semiclassical input
(Theorem~\ref{thm:hawking}).
\item Connects the cosmological hierarchy ($10^{120}$) to hidden geometry
via boundary dominance, with no semiclassical input.
\item Predicts $S = A/d$ for general observer dimension.
\item Performs the BTZ cross-check at $d = 3$ (Theorem~\ref{thm:btz}):
boundary dominance gives $S = A/3$, identification with standard
Bekenstein--Hawking entropy gives $G_d = d/4$ (Theorem~\ref{thm:Gd}),
and the cascade-derived Hawking temperature $T = r_h/(2\pi\ell^2)$
matches standard BTZ surface gravity exactly.
\end{itemize}

\textbf{Does not:}
\begin{itemize}
\item Compute metric-perturbation scattering amplitudes or demonstrate UV finiteness explicitly.
\item Derive the thermal spectrum and greybody factors from cascade geometry.
\item Compute entropy of Kerr or Reissner--Nordstr\"{o}m black holes.
\item Construct the explicit Hilbert space of metric superpositions or
compute decoherence rates.
\item Provide a \emph{second} cascade-internal route to
$G_d = d/4$ that does not go through the matching $S = A/d
\leftrightarrow A/(4G_d)$. Candidates (all cascade-native, none
Kaluza--Klein): the discrete cascade action of Paper~IVb~\cite{paper4b}
evaluated at layer~$d$; an independent sphere-area identity at $d$;
a direct calculation from the cascade's embedding data at that~$d$.
Theorem~\ref{thm:Gd} already derives $G_d$ cascade-internally from
boundary dominance; what is open is a \emph{corroborating} second route.
\item Derive the Page curve.
\end{itemize}

\section{Open Questions}\label{sec:open-questions-part2eq3}

\begin{enumerate}
\item\label{oq:state-metric-instantiation} \textbf{Instantiate the
state-to-metric map concretely.} Theorem~\ref{thm:metricstate} asserts
that different cascade states $|\Psi\rangle \in S^{d-1}$ produce
distinct 4D metrics $g_{\mu\nu}^{(\Psi)}$, under which superpositions
of states become superpositions of metrics
(Remark~\ref{rem:sp25-status}, Step~3). Paper~III's derivation of
$ds^2 = -dt^2 + a(t)^2\,d\sigma_3^2$ with $a(t) = \sqrt{1-t^2}$ yields
a single FRW background that does not depend on the state. What is
needed: a concrete pair of cascade states $|\Psi_1\rangle,
|\Psi_2\rangle$ together with a cascade-native construction exhibiting
distinct metrics $g^{(1)}_{\mu\nu} \neq g^{(2)}_{\mu\nu}$, via one of
two plausible routes---(a) state-dependent time direction in the
embedding $B^{d+1}$, producing distinct Lorentzian slicings of the
same underlying cascade geometry, or (b) state-dependent stress-energy
through the projection identity of Paper~III~\S8, sourcing distinct
solutions of Einstein's equation at $d=4$. Closing this question would
complete the ``gravity is already quantum'' claim by promoting the
metric-superposition picture from structural aspiration to concrete
construction. Related to the acknowledged gap on the explicit Hilbert
space of metric superpositions but logically distinct: this asks for
the underlying state-metric map, which is upstream of the Hilbert
space of metric states.
\item\label{oq:content-area-scale-invariance} \textbf{Derive the
scale invariance of the content--area map from a cascade action.}
Theorem~\ref{thm:bh} derives $V_d = \Omega_{d-1}/d$ on the unit ball
from boundary dominance, and identifies area with content via Paper~0
Corollary~3.2. The remaining step---that $S = A/d$ holds for any cascade
area $A$, not only the unit-sphere area $A = \Omega_{d-1}$---is asserted
via an additivity/counting principle (Remark~\ref{rem:sp23-status},
Step~3), with empirical support from Schwarzschild at $d=4$
(Theorem~\ref{thm:hawking}) and BTZ at $d=3$ (Theorem~\ref{thm:btz}).
What is needed: a derivation that forces the linear content--area map
from a cascade-native structural principle---e.g., a cascade action on
the lattice whose equilibrium content density per unit boundary area is
$V_d/\Omega_{d-1} = 1/d$ independent of the total area, or a
cascade-internal theorem that $S = A/d$ is the unique entropy
assignment consistent with boundary dominance plus the cascade's
content-equals-area primitive. Load-bearing on Bekenstein--Hawking,
Hawking temperature (Theorem~\ref{thm:hawking}), and the BTZ
cross-check (Theorem~\ref{thm:btz}).
\item \textbf{Hawking radiation spectrum: Araki--Woods route and
greybody factors.} Theorem~\ref{thm:hawking} derives the temperature
$T = 1/(8\pi M)$ from boundary dominance and Birkhoff's theorem, with
no semiclassical input. A cascade-internal route to the spectral
\emph{shape} exists but is not yet formally written up in the paper:
(i)~the cascade's slicing integrand $(1-x^2)^{d/2}$ is Gaussian
(Paper~II, Lemma~3.1); (ii)~the decoherence factor
$D = \exp(-D x^2/4R_{\rm eff}^2)$ is Gaussian (Paper~II,
Theorem~8.1); (iii)~by the Araki--Woods theorem, Gaussian
entanglement across a partition produces a reduced density matrix
with geometric eigenvalues $p_n = (1-r)r^n$, which is the
Bose--Einstein distribution. Numerically verified at $d=4$: the
cascade's Gaussian parameters $(a=15/4,\,b=7/4)$ give $r=0.2476$
matching the geometric series to four significant figures. Combined
with $T=1/(8\pi M)$ this predicts $n(\omega) = \Gamma(\omega)/
(\exp(\omega/T)-1)$. What remains open: (a)~formal theorem statement
and proof of the Araki--Woods step within the paper; (b)~greybody
factors $\Gamma(\omega)$, which require solving the
Regge--Wheeler/Zerilli wave equation on Schwarzschild---standard PDE
work downstream of the cascade's Einstein equation, but not yet
carried out; (c)~a geometric reading of the radiation
\emph{mechanism} (content leaking through asymptotic
compactification).
\item \textbf{Second cascade-internal route to $G_d = d/4$.}
Theorem~\ref{thm:Gd} derives $G_d$ cascade-internally from boundary
dominance plus matching to $A/(4G_d)$; Theorem~\ref{thm:btz} verifies
this is self-consistent with standard BTZ thermodynamics at $d=3$. What
remains open is a \emph{second} cascade-internal derivation of the same
$G_d = d/4$ relation via a distinct structural identity. This is
\emph{not} a Kaluza--Klein reduction (the cascade explicitly refuses
semiclassical reduction over compactified dimensions;
Paper~I~\S3.2~\cite{paper1}, Paper~II\,=\,III~\S5--6,
Paper~III~\S12~\cite{paper3}). Candidate cascade-native routes: the
proposed discrete cascade action (Paper~IVb~\cite{paper4b} Remark~4.6)
evaluated at layer~$d$, producing a variational derivation of $G_d$; an
independent sphere-area identity at layer~$d$ equal to $d/4$; or a
direct computation from the cascade's embedding data at dimension~$d$.
A second derivation would promote Theorem~\ref{thm:Gd} from
``structural matching'' to ``doubly-derived,'' tightening the BTZ
cross-check from self-consistency to independent corroboration.
\item \textbf{Metric-perturbation scattering.} The cascade predicts no
graviton particles (Paper~III, Section~12). Computing the effective
scattering amplitude of metric perturbations from the cascade's finite
tower and showing it is UV-finite would provide computational confirmation
of the structural finiteness.
\item \textbf{Decoherence of metric superpositions.} The Born rule applies
to metric measurements (Theorem~\ref{thm:metricstate}), but the rate at
which macroscopic metric superpositions decohere has not been computed.
\item \textbf{Holographic entropy bound.} Boundary dominance
($\Omega/V = d$) is structurally holographic. Whether the Bousso bound
follows from the cascade is open.
\item\label{oq:page-curve} \textbf{Page curve from cascade
evaporation.} A cascade-internal derivation proceeds in ten steps:
(1) total cascade state is pure (point on $S^{d-1}$, Paper~II);
(2) $S_{\rm BH} = 4\pi M^2$ from boundary dominance + Birkhoff;
(3) $T = 1/(8\pi M)$ (Theorem~\ref{thm:hawking});
(4) spectrum Planckian via the Araki--Woods route (Open
Question~3); (5) evaporation rate $dM/dt = -\alpha/M^2$ from
Stefan--Boltzmann + $A,T$; (6) $M(t) = (M_0^3 - 3\alpha t)^{1/3}$;
(7) $S_{\rm BH}(t) = 4\pi M(t)^2$;
(8) $S_{\rm rad}(t) = \min(S_{\rm emitted}(t), S_{\rm BH}(t))$
(Page's prescription, from purity at step~1); (9) Page time
$t_{\rm Page} = 0.646\, t_{\rm evap}$ where $S_{\rm rad}=S_{\rm BH}$;
(10) information recovery as $D\to 1$ when $A\to 0$. Key numbers:
$t_{\rm Page} = 0.646\, t_{\rm evap}$, maximum radiation entropy
$0.500\, S_0$, mass at Page time $M_0/\sqrt{2}$. Four open issues:
\begin{itemize}
\item \emph{Species coefficient.} The evaporation rate's
$\sigma$ depends on the number of emitting species, determined in
principle by the generation structure (Paper~IVa) but not yet
computed for BH evaporation.
\item \emph{Gaussian vs.\ random entanglement.} Page's
$S_{\rm rad}=\min(\cdot,\cdot)$ is exact for Haar-random states; the
cascade's entanglement is Gaussian. For $S_{\rm BH}\gg 1$ the two
agree to high accuracy, but cascade-specific corrections near the
Page time should be computable from the Gaussian eigenvalue structure.
\item \emph{Remnant question.} Whether the cascade's finite
213-layer tower permits complete evaporation ($A\to 0$) or forces a
Planck-scale remnant is not established; the remnant endpoint
modifies step~(10).
\item \emph{Backreaction.} The quasi-static evaporation of steps
(5)--(6) assumes flat-space radiation; cosmological-scale BHs need
FRW curvature corrections.
\end{itemize}
\item\label{oq:kerr-rn} \textbf{Kerr and charged black holes.} The
cascade's $S = A/d$ applies to any horizon (boundary dominance). For
Kerr (rotating) and Reissner--Nordstr\"{o}m (charged) BHs, $A$ depends
on $(M,J,Q)$; the no-hair theorem---a mathematical consequence of the
cascade's Einstein equation---makes these solutions unique. The
temperature and Page curve should follow by the same chain as
Schwarzschild with modified $A(M,J,Q)$; this has not been carried out
explicitly within the cascade framework.
\item\label{oq:decoherence-scrambling} \textbf{Cascade decoherence
deficit and BH scrambling time.} The Part~0 computes the
eigenvalue deficit $\epsilon = 6.11\%$ for the full 213-layer Gram
matrix---the cascade retains $93.9\%$ coherence across the full
descent. Whether this structural number connects to BH decoherence
properties (the scrambling time, the Page time for cascade-scale BHs,
or the Gaussian corrections of Open Question~8) is unexplored.
\end{enumerate}

\begin{thebibliography}{9}

\bibitem{paper0} RTAC, \href{\cascadebase/cascade-series-part0.pdf}{\textit{The Cascade Series~--- Part~0: The
Gamma-Function Structure of the Cascade}}, March 2026.

\bibitem{paper1} RTAC, \href{\cascadebase/cascade-series-part1.pdf}{\textit{The Cascade Series~--- Part~I: Scale Variance
from Orthogonality}}, March 2026.

\bibitem{paper2} RTAC, \href{\cascadebase/cascade-series-part2.pdf}{\textit{The Cascade Series~--- Part~II: Quantum
Mechanics from the Cascade}}, March 2026.

\bibitem{paper3} RTAC, \href{\cascadebase/cascade-series-part3.pdf}{\textit{The Cascade Series~--- Part~III: General
Relativity from the Cascade}}, March 2026.

\bibitem{paper4a} RTAC, \href{\cascadebase/cascade-series-part4a.pdf}{\textit{The Cascade Series~--- Part~IVa:
Standard Model from the Cascade --- Gauge Group, Symmetry Breaking,
and Three Generations}}, March 2026.

\bibitem{paper4b} RTAC, \href{\cascadebase/cascade-series-part4b.pdf}{\textit{The Cascade Series~--- Part~IVb:
Masses, Couplings, and Precision Predictions from the
Geometric-Topological Factorization}}, March 2026.

\bibitem{paper6} RTAC, \href{\cascadebase/cascade-series-part6.pdf}{\textit{The Cascade Series~--- Part~VI:
Time as Tower Growth}}, March 2026.

\bibitem{goroff} M.~H.~Goroff and A.~Sagnotti, ``The ultraviolet behaviour
of Einstein gravity,'' \textit{Nuclear Physics B} \textbf{266}, 709--736
(1986).

\bibitem{gleason} A.~M.~Gleason, ``Measures on the closed subspaces of a
Hilbert space,'' \textit{Journal of Mathematics and Mechanics} \textbf{6},
885--893 (1957).

\bibitem{lovelock} D.~Lovelock, ``The four-dimensionality of space and the
Einstein tensor,'' \textit{Journal of Mathematical Physics} \textbf{13},
874--876 (1972).

\bibitem{dewitt} B.~S.~DeWitt, ``Quantum theory of gravity.~I. The canonical
theory,'' \textit{Physical Review} \textbf{160}, 1113--1148 (1967).

\bibitem{hawking75} S.~W.~Hawking, ``Particle creation by black holes,''
\textit{Communications in Mathematical Physics} \textbf{43}, 199--220 (1975).

\bibitem{btz} M.~Ba\~{n}ados, C.~Teitelboim, and J.~Zanelli, ``Black Hole
in Three-Dimensional Spacetime,'' \textit{Physical Review Letters}
\textbf{69}, 1849--1851 (1992).

\bibitem{carlip-btz} S.~Carlip, ``The $(2+1)$-dimensional black hole,''
\textit{Classical and Quantum Gravity} \textbf{12}, 2853--2879 (1995).

\end{thebibliography}

\end{document}
